\section{An algorithmic example on two non-caterpillar binary trees}
\label{app:algorithmic-reverse-calculation}

This appendix evaluates the proof on two concrete non-caterpillar unrooted
binary phylogenetic trees with $N=25$ labelled leaves.  The same value of
$f_1$ is obtained in three algebraically different ways:
\begin{enumerate}[label=\textup{(\roman*)}]
\item sum the four nonnegative families in \eqref{eq:A30} over all principal charts;
\item evaluate the four defining levels $U-\frac12R-\frac12C+\frac14B$ directly;
\item enumerate the ordered-quadruple kernel on $X^4$.
\end{enumerate}
Before presenting the global totals, we display one principal chart, one completed
square, and its reversed bridge in full detail.

\subsection{The two concrete non-caterpillar trees}

Let $X=\{0,1,\ldots,24\}$.  Take the two trees represented by the following
Newick strings; line breaks inside each string are only for display.
\begin{verbatim}
T1 =
(((((((0,((19,20),23)),(14,24)),22),5),6),(9,10)),
(((((1,7),11),((2,((13,(16,18)),17)),12)),(8,15)),3),
(4,21));

T2 =
((((0,15),((2,7),24)),(10,20)),
(((((1,(((3,14),9),8)),19),11),5),18),
(((((4,(12,(13,21))),17),16),(22,23)),6));
\end{verbatim}
Each tree has $25$ leaves of degree one, $23$ internal vertices of degree
three, and $47$ edges.  To verify directly that neither graph is a
caterpillar, delete the leaves and consider the induced internal skeleton.
The numbers of skeleton vertices of degrees $1,2,3$ are
\[
\begin{array}{c|ccc}
\toprule
&\deg 1&\deg 2&\deg 3\\
\midrule
T_1&7&11&5\\
T_2&6&13&4\\
\bottomrule
\end{array}
\]
A binary caterpillar has a path as its internal skeleton and therefore has no
degree-three skeleton vertex.  The last column proves that both $T_1$ and
$T_2$ are non-caterpillar binary trees.

Each tree has
\[
 |E_2(T_1)|=|E_2(T_2)|=25-3=22
\]
quartet-relevant internal edges.  We use the following reproducible indexing.
For an internal edge $g$, let $S(g)$ be the smaller side of its split; it is
unique because $|X|=25$ is odd.  Order the internal edges first by $|S(g)|$
and then lexicographically by the increasing tuple of elements of $S(g)$.
Write the resulting orders as
\[
 E_2(T_1)=\{e_0,\ldots,e_{21}\},
 \qquad
 E_2(T_2)=\{f_0,\ldots,f_{21}\}.
\]
Thus the principal charts are
\[
 \chi_{pq}=(e_p,f_q),
 \qquad 0\le p,q\le21,
\]
and there are $22^2=484$ of them.  As in equation~\eqref{eq:A30}, put
\[
 A_\chi:=
 \frac14\mathcal S_{1,\chi}
 +\frac12\mathcal S_{2,\chi}
 +\frac12\mathcal S_{3,\chi}
 +\mathcal S_{4,\chi}.
\]
Then
\[
 f_1^{\mathrm{A.30}}(T_1,T_2)
 =\sum_{p=0}^{21}\sum_{q=0}^{21}A_{\chi_{pq}}.
\]

\subsection{One principal chart in full detail}

Consider
\[
 \chi_{7,6}=(e_7,f_6),
\]
where the smaller sides of the two principal splits are
\[
 S(e_7)=\{1,7,11\},
 \qquad
 S(f_6)=\{2,7,24\}.
\]
Choose the endpoint $v_-$ on the smaller side of each principal edge;
equivalently, the displayed orientation points away from the smaller side.
Endpoint atoms are ordered
by the edge order above, followed by the residual atoms.  The four row atoms
are
\[
\begin{array}{ccl}
\toprule
\text{index}&\text{sector}&\text{taxon block}\\
\midrule
0&a&R_0=\{1,7\}\\
1&b&R_1=\{2,12,13,16,17,18\}\\
2&b&R_2=X\setminus\{1,2,7,11,12,13,16,17,18\}\\
3&c&R_3=\{11\}\\
\bottomrule
\end{array}
\]
and the four column atoms are
\[
\begin{array}{ccl}
\toprule
\text{index}&\text{sector}&\text{taxon block}\\
\midrule
0&b&C_0=\{0,15\}\\
1&a&C_1=\{2,7\}\\
2&b&C_2=X\setminus\{0,2,7,15,24\}\\
3&c&C_3=\{24\}\\
\bottomrule
\end{array}
\]
Consequently,
\[
 r=\left(\frac2{25},\frac6{25},\frac{16}{25},\frac1{25}\right),
 \qquad
 c=\left(\frac2{25},\frac2{25},\frac45,\frac1{25}\right),
\]
\[
 (a_{i,k})=
 \begin{pmatrix}
 0&\frac1{25}&\frac1{25}&0\\
 0&\frac1{25}&\frac15&0\\
 \frac2{25}&0&\frac{13}{25}&\frac1{25}\\
 0&0&\frac1{25}&0
 \end{pmatrix},
\]
and
\[
 (x_{i,k})=
 \begin{pmatrix}
 -\frac4{625}&\frac{21}{625}&-\frac3{125}&-\frac2{625}\\
 -\frac{12}{625}&\frac{13}{625}&\frac1{125}&-\frac6{625}\\
 \frac{18}{625}&-\frac{32}{625}&\frac1{125}&\frac9{625}\\
 -\frac2{625}&-\frac2{625}&\frac1{125}&-\frac1{625}
 \end{pmatrix}.
\]
The row-index classes used below are
\[
\begin{aligned}
D_{\chi_{7,6}}
&=\{(0,1),(0,2),(1,0),(1,3),(2,0),(2,3)\},\\
I_{\chi_{7,6}}
&=\{(0,3),(1,2),(2,1),(3,0),(3,1),(3,2),(3,3)\}.
\end{aligned}
\]

\subsection{One completed square}

Take the term of $\mathcal S_{1,\chi_{7,6}}$ indexed by
\[
 (i,j,k)=(1,0,1).
\]
In the definition of $\mathcal S_{1,\chi_{7,6}}$, the two arguments of this
completed square are
\[
 X:=x_{\bar i,\bar k}+x_{i,\bar k}=-\frac{23}{625},
 \qquad
 Y:=x_{\bar i,k}+x_{i,k}=\frac{32}{625}.
\]
For this term,
\[
 c_k=\frac2{25},\qquad
 c_{\bar k}=\frac{22}{25},\qquad
 c_{k^+}=\frac3{25},
\]
and
\[
 a_{i,\bar k}=\frac15,\qquad
 a_{j,\bar k}=\frac1{25},\qquad
 a_{i,k}=a_{j,k}=\frac1{25}.
\]
The four completed-square terms in $\mathcal P_{1,0,1}(X,Y)$ are all
nonzero and equal, in order, to
\[
 \frac{49}{1739375000},\qquad
 \frac{11}{2343750},\qquad
 \frac{81}{976562500},\qquad
 \frac{81}{195312500}.
\]
Therefore
\[
\begin{aligned}
 \mathcal P_{1,0,1}\left(-\frac{23}{625},\frac{32}{625}\right)
 &={}
 \frac{49}{1739375000}
 +\frac{11}{2343750}
 +\frac{81}{976562500}
 +\frac{81}{195312500}\\
 &=\frac{85107103}{16306640625000}.
\end{aligned}
\]

\subsection{The four families in this chart}

Evaluating every permitted target in the same chart gives
\[
\begin{array}{lcc}
\toprule
\text{family}&\text{nonzero summands}&\text{exact sum}\\
\midrule
\mathcal S_{1,\chi_{7,6}}&16&1947947093/97839843750000\\
\mathcal S_{2,\chi_{7,6}}&16&366630403/48919921875000\\
\mathcal S_{3,\chi_{7,6}}&20&1912609/354492187500\\
\mathcal S_{4,\chi_{7,6}}&18&913296/29541015625\\
\bottomrule
\end{array}
\]
For one explicit $\Psi$-target term in $\mathcal S_{3,\chi_{7,6}}$, take
$(i,j,u)=(0,1,0)$.  Then
\[
 X:=x_{\bar i,\bar u}+x_{i,\bar u}=\frac3{625},
 \qquad
 \Psi^{\chi_{7,6}}_{0,1}(0)=\frac2{375},
\]
so that
\[
 \frac14\Psi^{\chi_{7,6}}_{0,1}(0)X^2
 =\frac14\cdot\frac2{375}\left(\frac3{625}\right)^2
 =\frac3{97656250}.
\]
Substitution into equation~\eqref{eq:A30} yields
\[
\begin{aligned}
A_{\chi_{7,6}}
={}&\frac14\frac{1947947093}{97839843750000}
+\frac12\frac{366630403}{48919921875000}\\
&+\frac12\frac{1912609}{354492187500}
+\frac{913296}{29541015625}\\
={}&\boxed{\frac{1841063809}{43484375000000}}.
\end{aligned}
\]

\subsection{The second square under bridge reversal}

The endpoint column atom $C_1=\{2,7\}$ corresponds to the edge $f_1$
adjacent to the principal edge $f_6$.  Hence the ordered endpoint chart is
\[
 \alpha=(f_6,f_1).
\]
Its three blocks are
\[
\begin{aligned}
K={}&\{2,7\},\\
H={}&\{24\},\\
L={}&\{0,1,3,4,5,6,8,9,10,11,12,13,14,15,16,17,18,19,20,21,22,23\}.
\end{aligned}
\]
Bridge reversal gives
\[
 \rho\alpha=(f_1,f_6),
 \qquad
 (K,H,L)\longmapsto(L,H,K).
\]
The reversed endpoint chart belongs to the principal chart
\[
 \chi_{7,1}=(e_7,f_1),
\]
and not to $\chi_{7,6}$.  The global sum in \eqref{eq:A30} therefore contains both
orientations:
\[
\begin{array}{lll}
\chi_{7,6},&\alpha=(f_6,f_1),&(K,H,L),\\
\chi_{7,1},&\rho\alpha=(f_1,f_6),&(L,H,K).
\end{array}
\]
In the reversed chart, the endpoint block $L$ has atom index $0$.  For the
displayed completed-square form, the four square coefficients pair as follows:
\[
\begin{array}{c|c|c}
\toprule
\text{second-square term in }\mathcal P^{\chi_{7,6}}_{1,0,1}
&\text{coefficient}
&\text{first-square term in }\mathcal P^{\chi_{7,1}}_{1,0,0}\\
\midrule
(A.22),\ Y^2&1/5750&(A.23),\ X^2\\
(A.23),\ Y^2&33/5000&(A.22),\ X^2\\
(A.24),\ Y^2&1/2500&(A.25),\ X^2\\
(A.25),\ Y^2&1/500&(A.24),\ X^2\\
\bottomrule
\end{array}
\]
Thus the sum of the four second-square coefficients in the left column and
the sum of their four first-square partners in the right column are both
$211/23000$.  This equality relates two different principal charts and is
precisely the
$(A.22)\leftrightarrow(A.23)$ and $(A.24)\leftrightarrow(A.25)$ coefficient
matching used by the symbolic proof.  Reversal does not identify the
subtotals of the two principal charts.

\subsection{First computation: the global four-family sum}

To display all $484$ principal-chart contributions compactly, group the row
edge indices into four disjoint ranges.  Put
\[
 D_*:=88349873334030000000000
\]
and
\[
\begin{aligned}
G_1&=509632929097785686741,&
G_2&=338603181900355216089,\\
G_3&=551304736571593597776,&
G_4&=697786193522003585794.
\end{aligned}
\]
The exact grouped sums are
\[
\begin{array}{ccc}
\toprule
\text{principal row-edge range}&\text{number of charts}&\text{contribution}\\
\midrule
0\le p\le4&110&G_1/D_*\\
5\le p\le10&132&G_2/D_*\\
11\le p\le15&110&G_3/D_*\\
16\le p\le21&132&G_4/D_*\\
\bottomrule
\end{array}
\]
These ranges contain all $484$ charts.  Consequently,
\[
\begin{aligned}
\sum_\chi A_\chi
&=\frac{G_1+G_2+G_3+G_4}{D_*}\\
&=\frac{2097327041091738086400}{88349873334030000000000}\\
&=\boxed{\frac{9273}{390625}}.
\end{aligned}
\]
Thus
\[
 \boxed{f_1^{\mathrm{A.30}}(T_1,T_2)=\frac{9273}{390625}}.
\]

\subsection[Second computation: the defining formula]
{Second computation: the defining formula for $f_1$}

For each pair $(e,f)$, retain the notation
\[
 U_{e,f}=T[\sigma_e,\tau_f],\quad
 R_{e,f}=T[\sigma_e^+,\tau_f],\quad
 C_{e,f}=T[\sigma_e,\tau_f^+],\quad
 B_{e,f}=T[\sigma_e^+,\tau_f^+].
\]
Direct exact enumeration of the four defining levels gives
\[
\begin{aligned}
\sum_{e\in E_2(T_1)}\sum_{f\in E_2(T_2)}U_{e,f}
&=\frac{32602}{78125}
                    =\frac{163010}{390625},&
\sum_{e\in E_2(T_1)}\sum_{f\in E_2(T_2)}R_{e,f}
&=\frac{48424}{78125}
                    =\frac{242120}{390625},\\
\sum_{e\in E_2(T_1)}\sum_{f\in E_2(T_2)}C_{e,f}
&=\frac{251002}{390625},&
\sum_{e\in E_2(T_1)}\sum_{f\in E_2(T_2)}B_{e,f}
&=\frac{371296}{390625}.
\end{aligned}
\]
Therefore
\[
\begin{aligned}
f_1^{\mathrm{direct}}
&=\sum_{e\in E_2(T_1)}\sum_{f\in E_2(T_2)}\left(
U_{e,f}-\frac12R_{e,f}-\frac12C_{e,f}+\frac14B_{e,f}
\right)\\
&=\frac{163010-121060-125501+92824}{390625}\\
&=\boxed{\frac{9273}{390625}}.
\end{aligned}
\]
Hence the two algebraically distinct accumulations agree exactly:
\[
\boxed{
f_1^{\mathrm{direct}}(T_1,T_2)
=f_1^{\mathrm{A.30}}(T_1,T_2)
=\frac{9273}{390625}}.
\]

This agreement holds only after the global sum is taken.  For the displayed
chart, write
$U_{76}:=U_{e_7,f_6}$ and analogously for $R_{76},C_{76},B_{76}$.  Direct
enumeration gives
\[
 U_{76}=\frac{256}{390625},\qquad
 R_{76}=\frac{502}{390625},\qquad
 C_{76}=\frac{524}{390625},\qquad
 B_{76}=\frac{989}{390625}.
\]
Thus its fixed-chart inclusion--exclusion subtotal is
\[
 U_{76}-\frac12R_{76}-\frac12C_{76}+\frac14B_{76}
 =-\frac{39}{1562500},
\]
whereas
\[
 A_{\chi_{7,6}}=\frac{1841063809}{43484375000000}>0.
\]
Bridge reversal reallocates individual second-square coefficients across
reversed principal charts.  Equation \eqref{eq:A30} preserves the global sum,
not each individual chart subtotal.

\subsection{Third computation: the ordered-quadruple kernel}

This calculation does not use the chart decomposition.  We use the signed
indexed sum $\mathfrak O_T$, antisymmetric side-pair function $\phi_o$, and
kernel $K_T$ from Section~\ref{sec:counting}.  For
the two trees in this appendix,
\[
 |\overrightarrow A_2(T_1)|=52,
 \qquad
 |\overrightarrow A_2(T_2)|=50,
\]
so $|\mathfrak O_{T_1}|=74$ and $|\mathfrak O_{T_2}|=72$, counting
indexed occurrences.

For $\boldsymbol q=(x,y,z,w)\in X^4$, put
\[
 c(x,y,z,w)
 :=\frac14K_{T_1}(x,y;z,w)K_{T_2}(x,y;z,w).
\]
The symbolic kernel identity \eqref{eq:ordered-quadruple-kernel} gives
\[
 N^4f_1(T_1,T_2)=\sum_{(x,y,z,w)\in X^4}c(x,y,z,w).
\]
Moreover, Lemma~\ref{lem:kernel-to-path-bridge} proves
$|c(x,y,z,w)|\le1/4$ for every tuple, including every repeated-label
pattern.
We evaluate this expression separately for all
\[
 25^4=390625
\]
ordered quadruples and classify each tuple by the number
$d=|\{x,y,z,w\}|$ of distinct labels.  The exact totals are
\[
\begin{array}{c|r|r}
\toprule
d&\text{number of ordered quadruples}&\sum c(x,y,z,w)\\
\midrule
1&25&0\\
2&4200&287\\
3&82800&9098\\
4&303600&-112\\
\midrule
\text{all}&390625&9273\\
\bottomrule
\end{array}
\]
For additional control of the repeated-label part, encode the equality pattern
of $(x,y,z,w)$ by assigning $0,1,2,3$ in order of first appearance.  The only
nonzero pattern totals are
\[
\begin{array}{c|c}
\toprule
\text{canonical equality pattern}&\text{exact contribution}\\
\midrule
0101, 0110&287/2\text{ each}\\
0102, 0112, 0120, 0121&4549/2\text{ each}\\
0123&-112\\
\bottomrule
\end{array}
\]
All other equality patterns contribute zero.  Hence the repeated-label
contribution is computed directly as
\[
 R_{\mathrm{rep}}=0+287+9098=9385,
\]
rather than being inferred by subtracting the distinct-label total from the
answer.

For the $d=4$ part, direct quartet enumeration gives
\[
 \Same(T_1,T_2)=4198,
 \qquad
 \Diff(T_1,T_2)=8452.
\]
These values sum to $\binom{25}{4}=12650$.  The symbolic $24$-ordering count
in the proof of Proposition~\ref{prop:counting-identity} gives subtotal $4$
for every common displayed quartet and $-2$ for every different displayed
quartet.  The program also checks these $24$-ordering subtotals separately
for every four-subset in this example.  Consequently,
\[
 4\Same-2\Diff
 =4\cdot4198-2\cdot8452
 =-112,
\]
which is exactly the $d=4$ row of the table.  Combining the distinct and
repeated parts gives
\[
 \sum_{(x,y,z,w)\in X^4}c(x,y,z,w)
 =-112+9385
 =9273.
\]
After normalization by $N^4=25^4=390625$,
\[
 f_1^{X^4\text{-enum}}(T_1,T_2)
 =\frac1{25^4}\sum_{(x,y,z,w)\in X^4}c(x,y,z,w)
 =\frac{9273}{390625}.
\]
Therefore all three calculations agree exactly:
\[
\boxed{
 f_1^{\mathrm{A.30}}(T_1,T_2)
 =f_1^{\mathrm{direct}}(T_1,T_2)
 =f_1^{X^4\text{-enum}}(T_1,T_2)
 =\frac{9273}{390625}}.
\]
Among the $390625$ ordered quadruples, $172324$ have nonzero contribution.
The largest observed absolute contribution of a single ordered quadruple is
$1/4$, attaining the universal bound from
Lemma~\ref{lem:kernel-to-path-bridge}.

All values in this appendix were recomputed in exact rational arithmetic by
the accompanying program \texttt{verify\_v048\_appendix\_c\_n25\_general.py}.
The program evaluates the global A.30 chart sum, compares every paired square
coefficient under bridge reversal, enumerates $X^4$, and verifies the stated
distinct-count and equality-pattern totals.  The A.30, defining-formula, and
$X^4$ calculations use algebraically different formulas on the same
deterministically generated tree data.  The program also compares the
canonical Newick serializations with the displayed strings, but the three
calculations do not use independently parsed tree implementations.

Accordingly, these finite calculations provide a three-formula consistency
check on shared parsed tree data.  They do not replace the row decomposition
with one column term fixed, the bridge-reversal argument, or the general
same-side block coefficient identity in
Appendix~\ref{app:local-expansion}; nor do they replace the full-domain
kernel-to-path proof in Section~\ref{sec:counting}.
